\documentclass[11pt]{beamer}
\usetheme{Madrid}
\usepackage{graphicx}
\usepackage{amsmath}
\usepackage{hyperref}
\usepackage{booktabs}
\usepackage{tikz}

\title[Polar Codes]{Shannon’s Theorem and Polar Codes}
\subtitle{From Theoretical Limits to Practical Encoding}
\author{Aly Assaf}
\institute{April 2025}
\date{}

\setbeamertemplate{footline}[frame number]
\setbeamertemplate{navigation symbols}{}

% Banner
\logo{\includegraphics[height=0.9cm]{example-image-a}} % Replace with your logo

%---------------------------------------------------
\begin{document}
\begin{frame}
  \titlepage
\end{frame}

%---------------------------------------------------
\section{Introduction}
\begin{frame}{Coding Theory History}
\begin{itemize}
    \item Shannon proved reliable communication is possible at rates below a channel’s capacity.
    \item But his proof was non-constructive—no method was given.
    \item Polar Codes (Arıkan, 2008) are the first provably capacity-achieving codes for symmetric binary-input memoryless channels.
    \item This talk introduces Coding Theory, Shannon’s Theorem, and the construction of Polar Codes.
\end{itemize}
\end{frame}


%---------------------------------------------------
\subsection{Example}
\begin{frame}{Coding Theory History}
\begin{itemize}
    \item Engineer's belief
\end{itemize}
\end{frame}
%---------------------------------------------------
\section{Foundations of Coding Theory}
\begin{frame}{The Coding Problem}
\begin{block}{Goal}
Design codes that:
\begin{itemize}
    \item Maximize transfer rate $R = \frac{k}{n}$
    \item Minimize error probability
\end{itemize}
\end{block}
\pause
\textbf{Key challenge:} Trade-off between efficiency ($R$) and reliability (distance $d$).
\end{frame}

\begin{frame}{Core Definitions}
\begin{itemize}
    \item \textbf{Code:} maps message $u \in \{0,1\}^k$ to codeword $x \in \{0,1\}^n$
    \item \textbf{Linear code:} if $Hx^T = 0$, then $x + y$ is also a codeword
    \item \textbf{Hamming weight:} number of 1s in a vector
    \item \textbf{Hamming distance:} $\text{dist}(x,y) = \text{wt}(x - y)$
    \item \textbf{Minimum distance:} smallest pairwise Hamming distance
\end{itemize}
\end{frame}

%---------------------------------------------------
\section{Shannon’s Theorem}
\begin{frame}{Information and Entropy}
\begin{itemize}
    \item \textbf{Information:} $h(x) = \log_2 \frac{1}{p(x)}$
    \item \textbf{Entropy:} average information content
    \[
    H(X) = \sum_{x} P(x) \log_2 \frac{1}{P(x)}
    \]
    \item \textbf{Mutual Information:} $I(X;Y) = H(X) - H(X|Y)$
\end{itemize}
\end{frame}

\begin{frame}{Shannon's Theorem}
\begin{block}{Theorem (1948)}
Reliable communication over a discrete memoryless channel is possible \\
if the rate $R < C$, where $C$ is the channel capacity.
\end{block}
\pause
\begin{itemize}
    \item For BSC(p), $C = 1 - h(p)$
    \item Converse: if $R > C$, error probability $\rightarrow 1$
    \item Shannon’s proof was existential—not constructive
\end{itemize}
\end{frame}

%---------------------------------------------------
\section{Polar Codes}
\begin{frame}{Why Polar Codes?}
\begin{itemize}
    \item Turbo and LDPC codes approached capacity only empirically
    \item Arıkan (2008) introduced the first provable, efficient method
    \item Polar codes achieve channel capacity with $O(n \log n)$ complexity
    \item Used in 5G; promising for 6G
\end{itemize}
\end{frame}

\begin{frame}{Key Idea: Polarization}
\begin{itemize}
    \item Start with two bits $U, V \sim \text{Bern}(p)$
    \item Encode as $(U + V, V)$ — reversible!
    \item Entropy shifts:
    \[
    H(U+V, V) = H(U,V) = 2h(p)
    \]
    \[
    H(V | U + V) = 2h(p) - h(p') < h(p)
    \]
\end{itemize}
\pause
\textbf{Some bits become more reliable, some less — this is polarization.}
\end{frame}

\begin{frame}{Recursive Construction}
\begin{itemize}
    \item Apply transformation recursively for $n = 2^t$ bits
    \item Output: bits polarized into
    \begin{itemize}
        \item Channels with entropy $\approx 0$ (good)
        \item Channels with entropy $\approx 1$ (bad)
    \end{itemize}
    \item Send only on good channels
\end{itemize}
\pause
\textbf{Preprocessing builds set $S$ of good indices.}
\end{frame}

%---------------------------------------------------
\section{Decoding and Performance}
\begin{frame}{Decoding Strategy}
\begin{itemize}
    \item Decoding is successive cancellation:
    \item Recover last bits first, then backtrack
    \item Efficient because structure is known: depth = $\log_2 n$
\end{itemize}
\pause
\textbf{Polar codes are practical:} Used in 5G NR standard (control channels).
\end{frame}

%---------------------------------------------------
\section{Conclusion}
\begin{frame}{Final Remarks}
\begin{itemize}
    \item Shannon showed what’s possible
    \item Arıkan showed how to achieve it
    \item Polar codes combine theory with efficient implementation
    \item They set a new standard for error-correcting codes
\end{itemize}
\end{frame}

%---------------------------------------------------
\begin{frame}[allowframebreaks]{References}
\begin{itemize}
    \item C. Shannon, “A Mathematical Theory of Communication,” 1948.
    \item E. Arıkan, “Channel Polarization: A Method for Constructing Capacity-Achieving Codes,” 2008.
    \item Additional sources as included in the original paper.
\end{itemize}
\end{frame}

\end{document}
